% ---------------------------------------------------------------------------
% Layer plans for multiblocks.
%
% A multiblock is built one layer at a time, from the bottom up, and that is how these draw it:
% a row of square grids, lowest on the left, each cell carrying a one-letter code for the block
% that goes there. It is how somebody actually builds -- lay a layer, step up, lay the next --
% rather than an isometric picture that looks better and is harder to follow with a stack of
% bricks in your hand.
%
% Drawn rather than screenshotted on purpose. A screenshot of a structure is wrong the moment
% the structure changes and nobody can tell by looking; a plan built from the same letters as
% its own key is at least internally consistent, and is a diff when it changes.
% ---------------------------------------------------------------------------

% ifthen supplies \ifthenelse, which the cell loops below use to test for an empty cell.
\usepackage{ifthen}
\usepackage{tikz}
\usetikzlibrary{positioning,fit,backgrounds}

\definecolor{mbcell}{RGB}{236,236,238}
\definecolor{mbedge}{RGB}{120,120,128}
\definecolor{mbmark}{RGB}{233,118,43}

% \mblayer{caption}{rows}
%
% `rows` is a semicolon-separated list of rows, top row first, each row a comma-separated list
% of cell letters. A dot means "nothing here"; a letter in braces is highlighted, which is how
% the block you strike with the hammer is marked.
%
% Rows are given top-first because that is how you read a plan on paper, and drawn with y
% decreasing so the first row ends up at the top.
\newcommand{\mblayer}[2]{%
  \begin{tikzpicture}[baseline=(current bounding box.north),
      cell/.style={draw=mbedge,fill=mbcell,minimum size=6.5mm,inner sep=0pt,
                   font=\sffamily\scriptsize},
      hot/.style={cell,fill=mbmark!30,draw=mbmark,very thick},
      gap/.style={minimum size=6.5mm,inner sep=0pt}]
    \foreach \row [count=\y] in {#2} {%
      \foreach \c [count=\x] in \row {%
        \ifthenelse{\equal{\c}{.}}
          {\node[gap] at (\x*6.7mm,-\y*6.7mm) {};}
          {\node[cell] at (\x*6.7mm,-\y*6.7mm) {\c};}%
      }%
    }
    \node[font=\sffamily\tiny,anchor=north] at (current bounding box.south) {#1};
  \end{tikzpicture}%
}

% The same, with one cell highlighted: \mblayerhot{caption}{rows}{col}{row}
\newcommand{\mblayerhot}[4]{%
  \begin{tikzpicture}[baseline=(current bounding box.north),
      cell/.style={draw=mbedge,fill=mbcell,minimum size=6.5mm,inner sep=0pt,
                   font=\sffamily\scriptsize},
      hot/.style={draw=mbmark,fill=mbmark!30,very thick,minimum size=6.5mm,inner sep=0pt,
                  font=\sffamily\bfseries\scriptsize},
      gap/.style={minimum size=6.5mm,inner sep=0pt}]
    \foreach \row [count=\y] in {#2} {%
      \foreach \c [count=\x] in \row {%
        \ifthenelse{\equal{\c}{.}}
          {\node[gap] at (\x*6.7mm,-\y*6.7mm) {};}
          {\ifthenelse{\x=#3 \AND \y=#4}
             {\node[hot] at (\x*6.7mm,-\y*6.7mm) {\c};}
             {\node[cell] at (\x*6.7mm,-\y*6.7mm) {\c};}}%
      }%
    }
    \node[font=\sffamily\tiny,anchor=north] at (current bounding box.south) {#1};
  \end{tikzpicture}%
}

% A row of layers with a shared key underneath.
%
% \begin{mbplan}{Coke Oven}{B = Coke Brick}
%   \mblayer{Layer 1}{B,B,B;B,B,B;B,B,B} \hfill ...
% \end{mbplan}
\newenvironment{mbplan}[2]
  {\par\vspace{6pt}\noindent\textbf{\sffamily\small #1}\\[2pt]%
   \def\mbplankey{#2}\noindent}
  {\\[2pt]{\footnotesize\itshape\mbplankey}\par\vspace{8pt}}
